\documentclass[11pt]{article}
\usepackage[margin=25mm]{geometry}
\usepackage{fontspec,unicode-math}
\setmainfont{Latin Modern Roman}
\setmathfont{Latin Modern Math}
\usepackage{amsmath,amssymb,mathtools}
\usepackage{xurl,hyperref}
\hypersetup{colorlinks=true,urlcolor=blue,linkcolor=blue}
\setcounter{secnumdepth}{0}
\setlength{\parindent}{0pt}
\setlength{\parskip}{5pt}
\setlength{\emergencystretch}{3em}
\newcommand{\sref}[2]{\hyperlink{source:#1:#2}{[S#2]}}
\newcommand{\cataloguescope}[1]{\paragraph{Source question.}#1}
\newcommand{\floor}[1]{\left\lfloor#1\right\rfloor}
\begin{document}
\section{Complexity of Dense Perfect Hypergraph Matching}\label{statement}

\subsection{Definitions and mathematical statement}

\paragraph{Shared notation.}Standard number systems are denoted by $\mathbb N,\mathbb Z,\mathbb Q,\mathbb R,\mathbb C$; index ranges, computational models and specialized notation follow the cited problem source and the definitions below.


A $k$-uniform hypergraph has edges that are $k$-element vertex subsets. A matching is a collection of pairwise disjoint edges. Its minimum $(k-1)$-codegree is the minimum number of edges containing a given $(k-1)$-subset.

A graph matching is a set of edges with pairwise disjoint endpoints, and it is perfect when every vertex is incident to one chosen edge. A weighted matching objective sums the chosen edge weights, unless an ordered objective is specified.
\paragraph{Statement recorded in the supplied source.}
For integers $k\geq2$ and $1/k\leq c<1/2$, determine the complexity of deciding whether a $k$-uniform hypergraph $H$ with $|V(H)|\equiv0\pmod k$ and $\delta_{k-1}(H)\geq cn$ has a matching of size $n/k$. \sref{statement}{2}
\paragraph{Source status.}
The supplied catalogue labels this record open. Its source-status metadata was checked against the supplied snapshot on 2026-10-11. This import does not independently establish current open status or mathematical correctness. The catalogue literature assessment is dated 2026-08-21.
\subsection{Short English statement}
Classify the difficulty of finding perfect matchings in uniform hypergraphs above the stated codegree threshold.
\subsection{Sources}
\begin{itemize}
\item[\hypertarget{source:statement:1}{S1}] Ulam AI and the UnsolvedMath Contributors, UnsolvedMath: A Curated Collection of Open Mathematics Problems, catalogue entry. Catalogue curation and metadata: CC BY 4.0. The supplied source question is retained; definitions and the English summary are editorial additions.
\url{https://huggingface.co/datasets/ulamai/UnsolvedMath}.
\item[\hypertarget{source:statement:2}{S2}] Combinatorics, Probability and Computing (2007), Oberwolfach Reports, DOI 10.4171/owr/2006/48, pp. 2875--2952
\url{https://doi.org/10.4171/owr/2006/48}.
\textit{Underlying Oberwolfach Reports material retains its source-specific licence, including CC BY-SA 4.0 where applicable; see THIRD\_PARTY\_NOTICES.md.}
\end{itemize}
\paragraph{Preparation and review.}
Original question attributed to the cited source and UnsolvedMath contributors. Codex (GPT-6; exact serving revision not exposed) assisted with definitions, English summaries, source inspection and formatting. No new human mathematical authorship or independent proof verification is asserted. Catalogue identity checked against the complete supplied registry and existing repository statements; known flagged repetitions and obvious identical conjectures excluded. Full mathematical-equivalence/literature review has not been performed. Statement preparation and final review followed contribution prompts C and D; automated checks verify structure and identity, not open status.
\label{end:statement}
\end{document}
